\documentclass[11pt]{article}
\usepackage[utf8]{inputenc}
\usepackage{amsmath,amssymb,amsthm}
\usepackage{graphicx}
\usepackage{hyperref}
\usepackage{booktabs}
\usepackage{xcolor}
\usepackage{geometry}
\usepackage{enumitem}
\usepackage{natbib}
\geometry{margin=1in}

% Theorem environments
\newtheorem{theorem}{Theorem}[section]
\newtheorem{lemma}[theorem]{Lemma}
\newtheorem{proposition}[theorem]{Proposition}
\newtheorem{corollary}[theorem]{Corollary}
\newtheorem{definition}[theorem]{Definition}
\newtheorem{remark}[theorem]{Remark}
\newtheorem{principle}{Principle}

\title{Limits of Self-Correction in LLMs:\\
\large An Information-Theoretic Analysis of Correlated Errors}

\author{Andrew Michael Brilliant\\
Independent Researcher\\
Sapporo, Japan\\
\texttt{a.brilliant@ieee.org}\\
ORCID: 0009-0004-8024-5442}

\date{January 2026}

\begin{document}

\maketitle

\begin{abstract}
We develop a diagnostic framework for evaluating when LLM self-evaluation can be trusted.
The framework's central results are: (1) under a shared-blind-spot modeling assumption with
joint conditional independence of evaluations given the shared failure structure, $k$ rounds
of self-critique provide information about correctness bounded by what the shared latent
failure variable $Z$ mediates---not by any independent channel---so that confidence
accumulated through repeated self-evaluation reflects the shared failure structure rather
than independently accumulated evidence; and (2) a selector provides positive evidence
about correctness if and only if it satisfies two independently measurable conditions: a
bounded false-acceptance rate and a true-acceptance rate exceeding that bound. These results
are conditional on explicit modeling assumptions. We also prove an information-theoretic
bound showing that when a shared latent failure structure mediates both generation and
evaluation errors, self-evaluation is bounded in what it can add---but we foreground this
result as scaffolding rather than the main contribution, since the latent variable is
unoperationalized and the bound is vacuous without independent constraints on it.

The diagnostic framework identifies what to measure to determine whether a deployed
system is in the failure regime, and what properties an external selector must have to
escape it. We describe design principles for an architecture motivated by this analysis.
Same-model context separation does not satisfy the formal independence conditions---it
is an engineering heuristic, not a theoretical solution---and we present it explicitly
as a practical starting point pending empirical validation.
\end{abstract}

\tableofcontents
\newpage

%%%%%%%%%%%%%%%%%%%%%%%%%%%%%%%%%%%%%%%%%%%%%%%%%%%%%%%%%%%%%%%%%%%%%%%%%%%%%%%
\section{Introduction}
%%%%%%%%%%%%%%%%%%%%%%%%%%%%%%%%%%%%%%%%%%%%%%%%%%%%%%%%%%%%%%%%%%%%%%%%%%%%%%%

Recent work demonstrates that large language models cannot reliably self-correct their
reasoning~\citep{huang2024selfcorrect}. \citet{tsui2025selfcorrection} provides further
evidence that this failure is not a knowledge deficit. LLMs can correct identical errors when
presented as external input but fail to correct those same errors in their own outputs, a
phenomenon termed the \textbf{Self-Correction Blind Spot}, measured at an average 64.5\%
failure rate across 14 models. We argue this reflects a structural property of certain
evaluation configurations: when generator and evaluator share failure modes, self-evaluation
provides weak evidence of correctness. A derivation may be elegant, internally consistent,
and convincingly presented, yet contain a subtle error that the model cannot detect in its
own output.

This failure is not primarily about compute or scale. It is about validation. Reliable
workflows require a mechanism that separates signal from noise, correct outputs from
plausible-but-wrong ones. The question is: what properties must a validation mechanism have
to be reliable?

\subsection{The Core Problem}

Consider a system that generates a hypothesis and then evaluates whether that hypothesis is
correct. Under what conditions does self-evaluation provide useful information?

We argue that the answer depends on error correlation. When the evaluator makes errors on
the same inputs where the generator makes errors, self-evaluation can be non-identifying:
agreement between generator and evaluator may provide weak evidence of correctness. This
echoes well-documented phenomena in human reasoning: confirmation bias, the curse of
knowledge, and why peer review and second opinions exist at all. The difference with LLMs
is that context can be deleted. A fresh instance has no memory of the reasoning it might
defend.

This is not a claim about any specific model's limitations. It is a structural property of
certain evaluation configurations. A single agent evaluating its own outputs faces elevated
risk of correlated error: sharing training data, inductive biases, and blind spots creates
the conditions under which the formal bounds of Section~\ref{sec:self_eval_fails} may
apply.\footnote{We say ``may apply'' deliberately. Whether these conditions hold for a
given deployed model and task is an empirical question, not a consequence of the formalism
alone.} \citet{tsui2025selfcorrection} offers empirical support for this structural claim.
Models that successfully correct errors in externally-presented solutions fail on their own
identical errors at substantially higher rates, suggesting that the knowledge to detect errors
exists but is not activated during self-evaluation.

\subsection{The Deep Context Challenge}

This problem becomes acute as context windows expand. Modern LLMs support large context
windows, enabling extended reasoning: complex derivations, multi-step analyses, and long
research sessions.

But large context is where correlated error accumulation may be most severe. Each reasoning
step inherits context from previous steps. Errors compound: an error at one step cascades
through subsequent steps, propagating through sequential
reasoning~\citep{tyen2024llmsstepbystep}. The longer the reasoning chain, the more
opportunity for self-reinforcing mistakes that become invisible within the context that
produced them.

Self-evaluation within a deep context may struggle to catch these errors, because the
evaluation inherits the same drift that produced the mistake. This creates a tension: the
contexts where rigorous evaluation matters most may be the contexts where self-evaluation is
least reliable.

Context-separated evaluation may help address this tension. By evaluating outputs in fresh
context, without the reasoning trace that produced them, we reduce the inheritance of
correlated error. The deeper the original context, the more valuable context separation may
become. We stress that whether fresh-context evaluation actually satisfies the conditions for
effective external selection (Section~\ref{sec:self_eval_fails}) is an empirical question;
we return to this in Section~\ref{sec:architecture}.

A complementary explanation comes from training data composition. \citet{tsui2025selfcorrection}
observes that standard training corpora consist overwhelmingly of human demonstrations of
correct reasoning, with few examples of a reasoner identifying and correcting their own
errors mid-stream. Models may thus lack the behavioral templates for self-correction. They
have learned to produce solutions, not to interrupt and revise them. This training gap
compounds the context-inheritance problem: models neither detect accumulated errors nor
possess trained patterns for correcting them.

\subsection{External Selection}

If correlated error is the problem, then evaluation in a modified or fresh context (where
the original reasoning trace is absent) may provide more independent signal. We call this
external selection. The key observation is that ``external'' refers to \textit{context}, not
necessarily to different models. A fresh instance of the same model, without access to the
reasoning chain that produced the candidate, may provide more independent critique because
the error-producing context is absent.

Multi-agent systems may succeed when they introduce external selection channels~\citep{du2023debate}:
formal proof checkers, executable tests, numerical invariants, independently-trained critics,
or even the same model under fresh context. The common element is reducing correlation
between generator and evaluator failure modes.

This motivates a practical architecture that separates:
\begin{itemize}
    \item \textbf{Generation}: High-entropy exploration of the hypothesis space
    \item \textbf{Selection}: Low-entropy evaluation under external constraints
    \item \textbf{Feedback}: Updating generation based on what survives selection
\end{itemize}

\subsection{Contributions}

\begin{enumerate}
    \item \textbf{Repeated self-critique bound} (Lemma~\ref{lem:repeated}): Under joint
    conditional independence of evaluations given the shared failure structure, $k$ rounds
    of self-critique provide information about correctness bounded by what the shared latent
    variable $Z$ mediates---not by any independent channel. Confidence accumulated through
    repeated self-evaluation therefore reflects the shared failure structure rather than
    independently accumulated evidence. This is the paper's primary formal result.
    \item \textbf{Sufficient conditions for positive evidence}
    (Proposition~\ref{prop:multi_agent}): A selector satisfying explicitly measurable
    bounds on false-acceptance and true-acceptance rates provides a quantifiable lower
    bound on evidence about correctness. These conditions are empirically testable using
    existing benchmark infrastructure~\citep{lin2024criticbench, kadavath2022ptrue}.
    \item \textbf{Information-theoretic scaffolding} (Theorem~\ref{thm:info_bound}): Under
    a shared-blind-spot modeling assumption, self-evaluation is bounded in what it can add
    about correctness. This result is true but its bound is vacuous without an independently
    constrained latent variable $Z$; we present it as framing for the above results rather
    than as an independent contribution.
    \item \textbf{The falsifiable $Z$ problem}: Identifying the construction of an
    independently-measurable blind-spot variable as the key open empirical problem the
    framework requires---and providing concrete candidate approaches
    \item \textbf{Design principles}: A framework for external selection motivated by the
    analysis, with explicit acknowledgment that same-model context separation is an
    engineering heuristic rather than a theoretical solution
\end{enumerate}

\subsection{Scope and Claims}

This paper makes a narrow conditional claim: \emph{if} a shared latent failure structure
mediates both generation errors and evaluation errors under joint conditional independence,
\emph{then} self-evaluation cannot add information about correctness beyond what that
structure already implies, and repeated self-critique accumulates confidence through the
shared failure structure rather than through independent evidence. We do not claim this
condition holds universally, nor that all self-evaluation fails. Whether the condition
holds for a given model and task is an empirical question, not a consequence of the
formalism.

The design principles we propose are \emph{motivated by} this analysis but are not
\emph{derived from} it: context separation using the same model does not satisfy the formal
independence conditions the theory identifies as necessary. We present the architecture as a
practical engineering starting point, explicitly disclaim theoretical justification for it,
and defer empirical validation to future work.

%%%%%%%%%%%%%%%%%%%%%%%%%%%%%%%%%%%%%%%%%%%%%%%%%%%%%%%%%%%%%%%%%%%%%%%%%%%%%%%
\section{Problem Formalization}
\label{sec:formalization}
%%%%%%%%%%%%%%%%%%%%%%%%%%%%%%%%%%%%%%%%%%%%%%%%%%%%%%%%%%%%%%%%%%%%%%%%%%%%%%%

\subsection{Basic Variables and Setup}

Let the input space be $\mathcal{X}$ and hypothesis space $\mathcal{H}$. For an input
$X \sim \mathcal{D}$, we define:

\begin{definition}[Generator and Selector]
Let $H^\star: \mathcal{X} \to \mathcal{H}$ denote the true (target) hypothesis function.
A \textbf{generator} produces candidate hypotheses via a conditional distribution
$G \sim p(g \mid x)$. A \textbf{selector} produces an evaluation score via
$S \sim p(s \mid x, g)$. The \textbf{acceptance decision} is a deterministic threshold:
\begin{equation}
A := \mathbb{I}\{S \geq \tau\}
\end{equation}
\end{definition}

\begin{definition}[Correctness Indicator]
We define the \textbf{correctness indicator}:
\begin{equation}
T := \mathbb{I}\{G = H^\star(X)\}
\end{equation}
Thus $T = 1$ when the generator output matches the true hypothesis, and $T = 0$ otherwise.
\end{definition}

\begin{definition}[Error Events]
The \textbf{generator error event} is $E_G := \{T = 0\}$. The \textbf{selector error
event} is:
\begin{equation}
E_S := \{A \neq T\}
\end{equation}
That is, $E_S$ occurs when the selector accepts an incorrect hypothesis ($A=1, T=0$) or
rejects a correct one ($A=0, T=1$).
\end{definition}

\begin{definition}[Conditional Error Coupling]
The \textbf{conditional error coupling} is:
\begin{equation}
\kappa := \Pr(E_S = 1 \mid E_G = 1) = \Pr(A = 1 \mid T = 0)
\end{equation}
This is the probability that the selector accepts given that the generator is wrong. Strong
coupling ($\kappa \approx 1$) indicates that the selector fails to detect generator errors.
\end{definition}

\subsection{Shared Blind Spots}

To model one mechanism by which errors might correlate, we introduce a latent variable
capturing shared failure modes.

\begin{definition}[Blind Spot Variable]
Let $Z \in \mathcal{Z}$ be a latent variable indexing input regions where a model family
systematically struggles. We say generator and selector \textbf{share blind spots} indexed
by $Z$ if there exists a set $\mathcal{Z}_{\mathrm{bad}} \subset \mathcal{Z}$ such that:
\begin{align}
\Pr(T = 0 \mid Z \in \mathcal{Z}_{\mathrm{bad}}) &\geq \alpha_G \gg \Pr(T = 0) \\
\Pr(A = 1 \mid T = 0, Z \in \mathcal{Z}_{\mathrm{bad}}) &\geq \alpha_S \gg \Pr(A = 1 \mid T = 0)
\end{align}
while both failure rates are low on $\mathcal{Z} \setminus \mathcal{Z}_{\mathrm{bad}}$.
\end{definition}

Shared blind spots arise naturally when generator and selector share training data,
architecture, or context. The variable $Z$ indexes specific failure modes (e.g., particular
reasoning patterns, domain gaps, or edge cases) where both components fail together.
For the formal results to have practical force, $Z$ should be given an independent,
a priori characterization---for example, as a fixed taxonomy of known failure modes for
a given model family---rather than being inferred post-hoc from observed joint failure.
\citet{erroratlas2026} provide one concrete instantiation: an empirically-derived taxonomy
of 17 error categories derived from analysis of 83 LLMs across 35 datasets, including
Logical/Reasoning Error, Calculation Error, Question Misinterpretation, and Incomplete
Reasoning. Each category is defined independently of any particular generator-selector pair
and can serve as an a priori $Z$ specification. The Self-Correction Blind Spot measured
by \citet{tsui2025selfcorrection}, in which models fail to correct their own errors despite
correcting identical errors in external input, is consistent with high $\alpha_S$ in the
self-evaluation setting. \citet{kim2025correlated} provide large-scale empirical evidence
for the error correlation assumption: across 350+ LLMs, same-family models agree on the
same wrong answer approximately 60\% of the time, and this correlation increases with model
capability --- directly measuring the high $\kappa$ regime the framework analyzes.

\begin{remark}[The modeling assumption $S \perp T \mid (G, Z)$]
\label{rem:ci_assumption}
Note that the selector $S$ is defined as $S \sim p(s \mid x, g)$: it has access to the
input $x$ as well as the generator output $g$. Since $T = \mathbb{I}\{G = H^\star(X)\}$
also depends on $x$, the selector could in principle be informative about $T$ by comparing
$g$ against the correct answer for the specific $x$---a channel independent of $Z$. The
conditional independence assumption $S \perp T \mid (G, Z)$ is a \emph{modeling assumption}
that captures the specific regime of interest: one where the selector's evaluation process
is itself corrupted by the same shared blind spots $Z$, so that the selector gains no
independent signal about correctness by accessing $x$ beyond what $Z$ already mediates.
This is precisely the self-evaluation setting, where a model evaluating its own output
shares the same failure modes on the same inputs. Selectors that can independently access
ground truth (e.g., a proof checker, test executor, or independently-trained model) violate
this assumption in the favorable direction, which is exactly why external selection works.
Readers should interpret all subsequent results as conditional on this assumption, which
characterizes a specific---empirically plausible---failure regime rather than a universal
property of self-evaluation.
\end{remark}

%%%%%%%%%%%%%%%%%%%%%%%%%%%%%%%%%%%%%%%%%%%%%%%%%%%%%%%%%%%%%%%%%%%%%%%%%%%%%%%
\section{When Self-Evaluation Fails}
\label{sec:self_eval_fails}
%%%%%%%%%%%%%%%%%%%%%%%%%%%%%%%%%%%%%%%%%%%%%%%%%%%%%%%%%%%%%%%%%%%%%%%%%%%%%%%

\subsection{Main Result}

The central quantity studied in this section is $I(T; S \mid G)$: the mutual information
between correctness $T$ and selector score $S$, conditional on the generator output $G$.
Conditioning on $G$ rather than on $(G, X)$ is a deliberate modeling choice with substantive
implications. Because $T = \mathbb{I}\{G = H^\star(X)\}$ depends on $X$, a selector with
access to $X$ could in principle be informative about $T$ through $X$ alone---a channel
independent of the shared blind-spot structure $Z$ (see Remark~\ref{rem:ci_assumption}).
Conditioning only on $G$ means this paper analyzes the regime where the selector's access
to $X$ does not provide an independent correctness signal beyond what $Z$ already mediates.
This is the self-evaluation regime of interest: a model evaluating its own output shares the
same failure modes on the same inputs, so the X-channel does not rescue evaluation quality.
Readers should interpret all results in this section as conditional on that regime. The
results do not apply to selectors that genuinely exploit $X$ independently---for example, a
proof checker that accesses the problem statement directly and verifies the answer against it.

The formal results below establish: when evaluator error is coupled with generator error via
a shared latent structure, self-evaluation becomes non-identifying---agreement provides
negligible evidence of correctness.

We formalize this through two theorems: an information-theoretic bound and an evidence bound.

\subsection{Information-Theoretic Formulation}

The central quantity is $I(T; S \mid G)$: the information that the selector provides about
correctness, given that we already observe the generator output.

\begin{theorem}[Information Bound via Shared Blind Spots]
\label{thm:info_bound}
Let $Z$ be a latent variable. Assume the conditional independence $S \perp T \mid (G, Z)$.
Then:
\begin{equation}
I(T; S \mid G) \leq I(T; Z \mid G)
\end{equation}
In particular, if $T \perp Z \mid G$, then $I(T; S \mid G) = 0$.
\end{theorem}

\begin{proof}
By the chain rule for conditional mutual information:
\begin{equation}
I(T; S, Z \mid G) = I(T; Z \mid G) + I(T; S \mid G, Z)
\end{equation}
Under $S \perp T \mid (G, Z)$, we have $I(T; S \mid G, Z) = 0$, hence:
\begin{equation}
I(T; S, Z \mid G) = I(T; Z \mid G)
\end{equation}
Also by the chain rule:
\begin{equation}
I(T; S, Z \mid G) = I(T; S \mid G) + I(T; Z \mid G, S) \geq I(T; S \mid G)
\end{equation}
since mutual information is nonnegative. Combining yields $I(T; S \mid G) \leq I(T; Z \mid G)$.
If additionally $T \perp Z \mid G$, then $I(T; Z \mid G) = 0$ and the result follows.
\qed
\end{proof}

\begin{remark}[On the informativeness of the bound]
\label{rem:vacuity}
The bound $I(T; S \mid G) \leq I(T; Z \mid G)$ is non-vacuous only when $Z$ is constrained
independently of $T$ and $S$. Because $Z$ is a latent variable, one could in principle
always choose $Z$ to encode all residual dependence between $S$ and $T$, rendering the
right-hand side large and the bound trivially satisfied. Theorem~\ref{thm:info_bound}
therefore establishes \emph{conditions} under which self-evaluation is weak---specifically,
when a pre-specified $Z$ characterizes the shared failure structure and $T \perp Z \mid G$
holds---rather than proving that self-evaluation is weak in all realistic systems. Practical
application of this theorem requires an a priori characterization of $Z$ (e.g., a fixed
taxonomy of known failure modes) that does not depend on the observed outcome.
\end{remark}

\begin{remark}[The bound does not imply weakness without a tight, small cap]
\label{rem:cap_tightness}
The inequality $I(T; S \mid G) \leq I(T; Z \mid G)$ establishes a structural constraint,
not a demonstration that self-evaluation is weak. The cap $I(T; Z \mid G)$ could be large
--- if $Z$ carries substantial information about correctness beyond what $G$ reveals, the
bound permits $S$ to be highly informative about $T$. The theorem implies weakness only
when two conditions hold jointly: (1) $Z$ is independently constrained so the bound is
non-vacuous, and (2) $I(T; Z \mid G)$ is shown to be small, so the cap is tight. The
paper does not independently establish (2) for real LLM systems. The theorem should
therefore be read as identifying a structural relationship --- information from $S$ is
$Z$-mediated --- rather than as proof that self-evaluation is empirically weak. The
Lemma~\ref{lem:repeated} and Proposition~\ref{prop:multi_agent} carry more direct
practical content because their conditions are expressible in directly measurable
probabilities.
\end{remark}

\textbf{Interpretation.} Under the modeling assumption of Remark~\ref{rem:ci_assumption} and
when $Z$ is appropriately constrained (see Remark~\ref{rem:vacuity}), the information the
selector provides about correctness is bounded by how much the blind spot variable $Z$
``knows'' about correctness beyond what the generator output already reveals. When $Z$ is a
pure nuisance variable (encoding only \textit{how} the system fails, not \textit{whether} it
fails), self-evaluation provides zero additional information \emph{within this regime}. This
should not be read as a universal claim: selectors that access ground truth independently of
$Z$ (e.g., by executing code or consulting a proof checker) fall outside the assumption's
scope and may provide substantial information about correctness.

\subsection{The Necessity of Falsifiable $Z$}
\label{sec:falsifiable_z}

Remark~\ref{rem:vacuity} notes that $Z$ can always be constructed post hoc to satisfy the
conditional independence assumption, making the bound in Theorem~\ref{thm:info_bound}
trivially true. This is not a defect to hide --- it is an important structural observation
about the framework's scope. Because any observed self-evaluation failure pattern admits a
latent-variable explanation under Theorem~\ref{thm:info_bound}, that theorem cannot by
itself predict or diagnose failure in real systems. It establishes what follows \emph{if}
the shared-blind-spot model applies --- it does not establish that the model applies.

This shifts the scientific burden precisely: the theorem's contribution is to identify
what would need to be measured to determine whether a system is in the failure regime.
Specifically, the theorem becomes empirically meaningful only when $Z$ is characterized
independently of the failure events being analyzed --- defined prior to observation,
estimable from separate data, and constrained to a fixed taxonomy rather than inferred
post hoc from joint failure.

In contrast, the paper's primary results --- the repeated self-critique bound
(Lemma~\ref{lem:repeated}) and the multi-agent advantage proposition
(Proposition~\ref{prop:multi_agent}) --- have more direct empirical content.
Lemma~\ref{lem:repeated} shows that $k$ evaluations under joint conditional independence
provide information about correctness bounded by what $Z$ mediates --- not by any
independent channel --- so accumulating evaluations cannot push information beyond the
$Z$-mediated ceiling; the key question for real systems is whether that joint independence
holds, which is testable given a specified $Z$.
Proposition~\ref{prop:multi_agent} gives conditions a selector must satisfy to provide
positive evidence, expressed in terms of directly measurable probabilities $\delta$ and
$\epsilon$.

For $Z$ to ground Theorem~\ref{thm:info_bound} in real systems, candidate approaches
include: behavioral clustering of error types across prompts, probing studies of model
internals, and difficulty-stratified benchmarks designed to identify systematic failure
regions for a given model family. \citet{erroratlas2026}'s 17-category ErrorAtlas taxonomy
demonstrates that such a priori specifications are achievable: the taxonomy is derived from
held-out model behavior and defines failure categories independently of any particular
generator-selector relationship.

For the quantities in Proposition~\ref{prop:multi_agent} and Theorem~\ref{thm:evidence_bound},
measurement infrastructure already exists. \citet{lin2024criticbench} introduce CriticBench,
which evaluates the same LLM on generation, binary critique (correctness classification),
and correction for identical problems --- directly yielding $\Pr(A=1 \mid T=0)$ and
$\Pr(A=1 \mid T=1)$ from the critique task across 17 models and 5 reasoning domains.
\citet{kadavath2022ptrue} introduce the P(True) framework, measuring model-generated
self-assessments of correctness against known ground truth, reporting calibration curves
that are precisely these conditional probabilities. These existing protocols provide a
concrete experimental pathway to testing whether a given deployed system satisfies the
proposition's conditions, without requiring new benchmark construction.

Constructing a $Z$ that is both independently-specified and informative enough to make
Theorem~\ref{thm:info_bound} non-vacuous remains the key open empirical problem, but the
tools for doing so are available in the existing literature.

\begin{remark}[Ontological vs.\ epistemological status of $Z$]
\label{rem:z_ontology}
The framework requires care about two distinct claims involving $Z$. The \emph{ontological}
claim is that some $Z$ exists making the conditional independence $S \perp T \mid (G, Z)$
hold approximately---this is mathematically guaranteed by construction and carries no
empirical content on its own. The \emph{epistemological} claim is that such a $Z$ can be
independently specified, measured, and constrained prior to observing joint failure. Only
the epistemological claim gives the framework practical diagnostic power. The theorems
establish what follows \emph{if} an appropriate $Z$ is known; they do not establish that
any particular $Z$ is knowable. Practical application requires bridging this gap, which is
why we identify the falsifiable $Z$ construction as the key open problem rather than
treating it as solved by the theoretical existence result.
\end{remark}

\begin{lemma}[Post-Processing Cannot Increase Evidence]
\label{lem:post_processing}
For any deterministic acceptance rule $A = \mathsf{acc}(S)$:
\begin{equation}
I(T; A \mid G) \leq I(T; S \mid G)
\end{equation}
\end{lemma}

This follows directly from the data processing inequality. The acceptance decision cannot
contain more information than the selector score from which it derives.

\subsection{Evidence Bound Formulation}

\begin{theorem}[Bounded Evidence from Acceptance]
\label{thm:evidence_bound}
Assume the selector has high false acceptance rate:
\begin{equation}
\Pr(A = 1 \mid T = 0) \geq 1 - \varepsilon
\end{equation}
and $\Pr(A = 1 \mid T = 1) \leq 1$. Then the log-likelihood ratio contributed by observing
$A = 1$ satisfies:
\begin{equation}
\log_2 \frac{\Pr(A = 1 \mid T = 1)}{\Pr(A = 1 \mid T = 0)} \leq \log_2 \frac{1}{1 - \varepsilon}
\end{equation}
\end{theorem}

\begin{proof}
We have $\Pr(A = 1 \mid T = 1) \leq 1$ and $\Pr(A = 1 \mid T = 0) \geq 1 - \varepsilon$,
so:
\begin{equation}
\frac{\Pr(A = 1 \mid T = 1)}{\Pr(A = 1 \mid T = 0)} \leq \frac{1}{1 - \varepsilon}
\end{equation}
Taking $\log_2$ gives the result. \qed
\end{proof}

\begin{corollary}[Degenerate Evidence for Small $\varepsilon$]
For small $\varepsilon$:
\begin{equation}
\log_2 \frac{1}{1 - \varepsilon} \approx \frac{\varepsilon}{\ln 2} \approx 1.44\varepsilon
\text{ bits}
\end{equation}
To illustrate the mathematical behavior of the bound: for $\varepsilon = 0.01$ (selector
accepts 99\% of incorrect hypotheses), acceptance provides at most $0.014$ bits of evidence.
For $\varepsilon = 0.35$ (a more moderate regime), the bound rises to approximately $0.63$
bits. These are mathematical illustrations of how the bound scales with $\varepsilon$, not
empirical estimates; determining which regime applies to a given deployed system requires
direct measurement of acceptance probabilities.
\end{corollary}

\textbf{Important caveat.} Theorem~\ref{thm:evidence_bound} provides an \emph{upper} bound
on evidence under the assumption of high false acceptance. It does not establish that
acceptance is negligible in specific realistic settings unless the false acceptance rate is
empirically near 1. The bound applies specifically to the regime where $\Pr(A=1 \mid T=0)
\geq 1-\varepsilon$ for small $\varepsilon$; outside this regime, acceptance may convey
substantial evidence. Many successful LLM applications violate the preconditions by
incorporating external feedback channels (execution, formal verification, retrieval). The
negative results apply specifically to the regime where: (1) systematic generator failures
exist with nontrivial probability, and (2) the selector shares the generator's blind spots.
We do not claim all self-evaluation fails, only that it may fail under these conditions.

\textbf{Empirical context (non-load-bearing).} The theorem's precondition---that false
acceptance is elevated in self-evaluation settings---is not established by the formal
results but is contextually motivated by available empirical evidence. We note explicitly
that the mapping is loose: \citet{tsui2025selfcorrection}'s 64.5\% Self-Correction Blind
Spot measures ``failure to self-correct,'' which is not operationally identical to the
formal acceptance event $A=1$ (it could reflect abstention, switching to a different wrong
answer, or other behaviors). Direct substitution into Theorem~\ref{thm:evidence_bound}
would overstate precision. We cite this evidence only to note that the high-false-acceptance
regime the theorem analyzes is not implausible, not to establish that any real system is
in that regime. Precise empirical measurement of $\Pr(A=1 \mid T=0)$ requires controlled
experiments with explicit threshold decisions, which we leave to future work. The formal
results stand independent of whether Tsui's findings map onto the theorem's preconditions.

\subsection{The Confidence Amplification Problem}

Worse than providing no information, correlated self-evaluation can amplify confidence in
errors.

\begin{lemma}[Repeated Self-Critique Bound]
\label{lem:repeated}
Consider $k$ selector outputs $S_1, \ldots, S_k$ with acceptance decisions
$A_i = \mathbb{I}\{S_i \geq \tau\}$. Assume $(S_1, \ldots, S_k) \perp T \mid (G, Z)$
jointly---which holds, for example, when the $S_i$ are mutually conditionally independent
given $(G, Z)$ and each satisfies $S_i \perp T \mid (G, Z)$. Then:
\begin{equation}
I(T; A_{1:k} \mid G) \leq I(T; Z \mid G)
\end{equation}
That is, $k$ critiques provide no more information about correctness than the single blind
spot variable $Z$.
\end{lemma}

\begin{proof}
By the joint conditional independence assumption, $A_{1:k} \perp T \mid (G, Z)$. The chain
rule gives:
\begin{equation}
I(T; A_{1:k} \mid G) \leq I(T; Z \mid G) + I(T; A_{1:k} \mid G, Z) = I(T; Z \mid G)
\end{equation}
since $I(T; A_{1:k} \mid G, Z) = 0$ by the joint conditional independence. \qed
\end{proof}

\begin{remark}
The joint independence assumption in Lemma~\ref{lem:repeated} is stronger than assuming
each $S_i \perp T \mid (G,Z)$ individually; individual conditional independence does not in
general imply joint conditional independence. The joint assumption holds naturally when
selector evaluations are drawn independently from the same conditional distribution
$p(s \mid g, z)$, as in the case of $k$ independent fresh-context evaluations of the same
output under the same latent blind spot.
\end{remark}

\begin{proposition}[Confidence Amplification]
Under strongly coupled error, repeated self-evaluation that produces consistent acceptance
increases subjective confidence while providing no objective evidence.
\end{proposition}

\begin{proof}
If a system evaluates its hypothesis $k$ times and accepts each time, a naive Bayesian
update treats these as independent evidence:
\begin{equation}
\Pr(T = 1 \mid A_1, \ldots, A_k) \propto \Pr(T = 1) \prod_{i=1}^{k}
\frac{\Pr(A_i \mid T = 1)}{\Pr(A_i \mid T = 0)}
\end{equation}

Under strong error coupling, however, $A_1, \ldots, A_k$ are not independent conditional on
$(T, Z)$. If $T = 0$ and $Z \in \mathcal{Z}_{\mathrm{bad}}$, all evaluations fail together:
\begin{equation}
\Pr(A_1 = \cdots = A_k = 1 \mid T = 0, Z \in \mathcal{Z}_{\mathrm{bad}}) \approx 1
\end{equation}

By Lemma~\ref{lem:repeated} (under the joint independence assumption), the $k$ acceptances
collectively provide information about correctness bounded by what $Z$ mediates --- they
cannot push beyond the $Z$-mediated ceiling regardless of how many accumulate. The apparent
evidence of $k$ consistent acceptances is $Z$-mediated information, not $k$ independent
pieces of evidence. But the subjective experience is $k$ ``confirmations,'' creating false
confidence.
\qed
\end{proof}

This may contribute to a failure mode observed in extended LLM reasoning: increasing
confidence in coherent, well-argued, wrong conclusions.

\subsection{When External Selection Works}

The results above identify conditions under which self-evaluation may provide weak evidence.
The contrapositive suggests when external selection may help:

\begin{corollary}[External Selection Criterion]
\label{cor:external_selection}
Evaluation provides substantial information about correctness when:
\begin{enumerate}
\item The selector accesses information not contained in $(G, Z)$, breaking the conditional
independence $S \perp T \mid (G, Z)$
\item The selector's blind spots $\mathcal{Z}'_{\mathrm{bad}}$ have low overlap with the
generator's $\mathcal{Z}_{\mathrm{bad}}$
\item The false acceptance rate satisfies $\Pr(A = 1 \mid T = 0) \ll 1$
\end{enumerate}
\end{corollary}

External selection channels that satisfy these criteria include: formal verification
(accesses mathematical ground truth), executable tests (accesses computational ground truth),
different model families (different $\mathcal{Z}_{\mathrm{bad}}$), and fresh-context
evaluation (conjectured to partially reduce shared context-level blind spots; see
Section~\ref{sec:architecture} for caveats). The empirical results of
\citet{tsui2025selfcorrection} are consistent with the general picture: the ``Wait''
intervention, which injects a correction marker into the model's reasoning trace, reduced
the blind spot rate by 89.3\% without changing model weights. This is consistent with
context perturbation reducing error correlation, though the cited evidence establishes
performance improvement, not which formal conditional independence relations changed.
The mechanism remains a conjecture.

\subsection{Mechanistic Note: A Predictive Interpretation}

The information-theoretic results above establish \textit{that} self-evaluation can fail
under correlated error. This section offers one interpretation of \textit{why} such
correlation may be elevated in language models. This mechanistic account is consistent with
the formal analysis but does not depend on it; the information-theoretic bounds hold
regardless of the underlying mechanism.

\citet{bender2021parrots} characterize language models as systems that ``stitch together
sequences of linguistic forms... according to probabilistic information about how they
combine, but without any reference to meaning.'' Under this framing, when asked to evaluate
a hypothesis, a language model predicts what evaluative text would likely follow the prompt,
given its training distribution. This prediction inherits whatever patterns characterize
human evaluation behavior in that distribution: prestige deference, format heuristics,
social smoothing, and narrative continuation.

Alignment training (RLHF) shifts \textit{which} human behavior is predicted but may not
change the underlying operation. \citet{sharma2024sycophancy} demonstrate that RLHF-trained
models systematically exhibit sycophancy (responses matching user beliefs over truthful ones)
and that human preference judgments favor sycophantic responses, creating a training signal
toward agreement.

\paragraph{A note on optimization targets.} Standard system prompts optimize for
``helpful assistant,'' not ``rigorous evaluator'' or ``truth-seeker.''
\citet{zheng2023helpful} systematically evaluated social roles in system prompts and found
that the ``helpful assistant'' framing is nearly universal in commercial deployments, yet
produces measurably different behavior than alternative framings. These are not equivalent
objectives.

\paragraph{A hypothesis on format consistency.} One observable phenomenon follows from this
analysis: \textit{format consistency}. We hypothesize that submitting manuscripts of widely
varying quality to major language models with neutral prompts will yield near-uniform
response format---balanced positive and negative points regardless of input quality. This
hypothesis is stated in falsifiable form to invite experimental follow-up; we do not treat
it as an established result. Readers who wish to test it empirically are encouraged to do so
with appropriate controls for prompt variation and model version.

\paragraph{User control through prompting.} The behaviors described above are defaults, not
constraints. Extensive research demonstrates that prompt design substantially affects model
behavior, with performance differences of up to 76 percentage points from formatting changes
alone~\citep{sclar2024formatting}. Role prompting shifts reasoning performance dramatically;
\citet{kong2023roleplay} report accuracy improvements from 53.5\% to 63.8\% on mathematical
reasoning simply by changing the prompt framing. \citet{zhuo2024prosa} introduce sensitivity
metrics showing that prompt variations produce substantial and measurable behavioral shifts.
If you prompt a model with explicit evaluation criteria (``identify all flaws,'' ``be
maximally critical,'' ``act as a hostile reviewer seeking reasons to reject''), it will shift
toward that behavior.

\paragraph{Implication for context separation.} When a model generates output under a
``helpful assistant'' framing and then evaluates that output under the same framing, error
correlation risk is structurally elevated: both generation and evaluation are shaped by the
same training signal, which satisfies a necessary (though not sufficient) condition for the
shared blind spot model of Section~\ref{sec:formalization} to apply. Context separation may
help because a fresh context with explicit critic framing resets the prediction target,
potentially reducing the shared context-level component of $Z$.

Importantly, the system prompt in commercial deployments (ChatGPT, Claude, Gemini) is not
user-editable. User instructions are layered on top of this hidden foundation. A prompt like
``evaluate this critically'' operates atop ``be helpful,'' not instead of it. Agentic
frameworks built on these APIs inherit the same constraint.

%%%%%%%%%%%%%%%%%%%%%%%%%%%%%%%%%%%%%%%%%%%%%%%%%%%%%%%%%%%%%%%%%%%%%%%%%%%%%%%
\section{Selection Pressure Across Domains}
%%%%%%%%%%%%%%%%%%%%%%%%%%%%%%%%%%%%%%%%%%%%%%%%%%%%%%%%%%%%%%%%%%%%%%%%%%%%%%%

This observation is not specific to AI systems. Selection pressure (a mechanism that
determines which configurations persist) appears across biological, physical, and scientific
domains. We present these parallels as motivation for treating external selection as a
general pattern rather than a domain-specific observation.

\subsection{Self-Reference Limitations}

G\"{o}del's incompleteness theorems establish that sufficiently powerful formal systems
cannot prove their own consistency~\citep{godel1931}. The structural parallel is suggestive:
self-reference creates blind spots. A system that generates claims cannot fully validate
those claims using only internal resources. \citet{tsui2025selfcorrection} draws an
analogous parallel to cognitive science, connecting the LLM self-correction blind spot to
the bias blind spot documented by \citet{pronin2002bias}. Humans reliably identify cognitive
biases in others while failing to detect the same biases in themselves.

This mirrors a familiar experience in software engineering: developers cannot effectively QA
their own code. The problem is not laziness or lack of intelligence. It is that the developer
knows how the code is \textit{supposed} to work and cannot clear that context when testing.
The same cognitive patterns that produced the bug prevent recognizing it as a bug.

\subsection{Selection Pressure Across Domains}

Selection pressure provides the external criterion that distinguishes signal from noise. We
observe this structure across engineering, scientific, and reasoning domains:

\begin{table}[ht]
\centering
\caption{Selection Pressure Across Domains}
\small
\begin{tabular}{llll}
\toprule
\textbf{System} & \textbf{Perturbation} & \textbf{Selection Criterion} & \textbf{Amplification} \\
\midrule
\multicolumn{4}{l}{\textit{Engineering}} \\
Fuzzer & Random mutation & No crash & Bug-free path found \\
Monte Carlo & Stochastic proposal & Satisfies constraints & Solution region mapped \\
Genetic alg. & Crossover/mutation & Fitness improves & Optimized design \\
\midrule
\multicolumn{4}{l}{\textit{Scientific}} \\
Hypothesis & Conjecture & Persists under test & Theory in literature \\
\midrule
\multicolumn{4}{l}{\textit{Human Reasoning}} \\
Draft & Initial attempt & Survives fresh review & Revised manuscript \\
\midrule
\multicolumn{4}{l}{\textit{LLM (proposed)}} \\
High temp. & Random token & Survives fresh context & Validated output \\
\bottomrule
\end{tabular}
\end{table}

The pattern is consistent: perturbation generates variation, selection determines what
persists, and surviving configurations amplify.

We argue this pattern more closely reflects how human reasoning actually works. A theory is
rarely written in a single session. It is returned to with fresh eyes the next morning,
reviewed by colleagues with uncorrelated blind spots, revised after a week away from the
problem. Each return provides external selection: the researcher encounters only the output,
not the reasoning trace that produced it.

Consider the contrast with extended chain-of-thought in a single context. The model
generates a draft, evaluates it, and declares it ready, all while anchored to the reasoning
that produced the draft. In our experience, a manuscript declared ``ready'' in a long context
session will be identified as flawed when pasted into a fresh context. This observation
motivated the theoretical framework but does not constitute validation of it
(Section~\ref{sec:observations}).

The engineering cases in Table~1 show this pattern is already standard practice. Fuzzers
generate random inputs; programs that survive without crashing have demonstrated robustness.
Monte Carlo methods propose random configurations; those satisfying constraints map the
viable solution space. In each case, perturbation without selection produces nothing;
perturbation with external selection produces progress.

\subsection{External Selection in Practice: Formal Verification}

Recent work in theorem proving illustrates external selection concretely. The Prover Agent
framework~\citep{baba2025prover} coordinates an informal reasoning LLM with the Lean proof
assistant, where Lean provides external verification. Using relatively small language models,
the system achieved 88.1\% accuracy on the MiniF2F benchmark, outperforming approaches
using larger models without external verification.

The mechanism aligns with our analysis. The LLM generates candidate proofs (high-entropy
generation). Lean verifies whether the proof compiles, a criterion external to the LLM's
training and biases (external selection). The external selection channel reduces correlation
between generator error and evaluator error, because Lean's verification depends on
mathematical truth, not on patterns in training data.

\subsection{Persistence as the Criterion for Reliability}

There is a practical equivalence here that bears stating plainly: in any workflow,
\textit{what we treat as reliable} is \textit{what survives independent checks}. These are
not two separate properties; they are one property described from two directions.

In scientific practice, we can only build on what we can measure and replicate. What persists
under repeated, independent measurement is what we call ``established.'' This is not a
limitation of our methods; it is the operational definition of reliability.

The context-separated architecture is motivated by this principle. One context's output
becomes another context's input for critique, removing the generation trace from the
evaluation context. We stress again that this engineering heuristic is not theoretically
equivalent to the external selection the formal analysis identifies as necessary: same-model
context separation does not break parameter-level correlations. The analogy to independent
checks is motivating rather than exact.

%%%%%%%%%%%%%%%%%%%%%%%%%%%%%%%%%%%%%%%%%%%%%%%%%%%%%%%%%%%%%%%%%%%%%%%%%%%%%%%
\section{Multi-Agent Verification}
%%%%%%%%%%%%%%%%%%%%%%%%%%%%%%%%%%%%%%%%%%%%%%%%%%%%%%%%%%%%%%%%%%%%%%%%%%%%%%%

\emph{Note on scope: This section discusses multi-agent verification patterns motivated by
the formal analysis. The connections drawn are conceptual---the theory identifies what
properties an effective external selector must have; the discussion below catalogs
mechanisms that plausibly satisfy those properties. None of the claims in this section are
formally derived from the theorems.}

If self-evaluation fails under correlated error, how can multi-agent systems succeed?

\subsection{Breaking Correlation}

Multi-agent verification helps when it introduces selectors whose error is less correlated
with the generator.

\begin{definition}[External criterion]
An external criterion is a selection mechanism depending on information not fully controlled
by the generator.
\end{definition}

\begin{proposition}[Sufficient conditions for positive evidence]
\label{prop:multi_agent}
Suppose a selector satisfies $\Pr(A = 1 \mid T = 0) \leq 1 - \epsilon$ and
$\Pr(A = 1 \mid T = 1) \geq \delta$ for some $\delta > 1 - \epsilon > 0$. Then acceptance
by that selector provides
\begin{equation}
\log_2 \frac{\delta}{1 - \epsilon} > 0
\end{equation}
bits of evidence about correctness.
\end{proposition}

\begin{proof}
The log-likelihood ratio of $A=1$ is
$\log_2 \frac{\Pr(A=1 \mid T=1)}{\Pr(A=1 \mid T=0)} \geq \log_2 \frac{\delta}{1-\epsilon}$,
which is positive whenever $\delta > 1 - \epsilon$. \qed
\end{proof}

\begin{remark}
The proposition gives \emph{sufficient} conditions under which acceptance provides positive
evidence, not necessary and sufficient conditions. Positive evidence requires only
$\Pr(A=1 \mid T=1) > \Pr(A=1 \mid T=0)$; the $(\delta, \epsilon)$ parameterization makes
the bound explicit and measurable. A selector meeting these conditions provides at least
$\log_2(\delta/(1-\epsilon))$ bits; the actual evidence may be larger. The proposition
additionally requires the selector to accept correct hypotheses at rate $\geq \delta >
1-\epsilon$, since a trivially rejecting selector would satisfy the false-acceptance bound
without being useful.
\end{remark}

The key insight is diversity of failure modes. A selector trained on different data, using
different architecture, or implementing formal verification will fail on different inputs
than the generator.

\subsection{External Selection Channels}

Effective external selection channels include:

\paragraph{Formal verification.} Proof assistants (Lean, Coq, Isabelle) and type checkers
provide selection under mathematical ground truth. A proof that compiles is verified by
mathematics itself, not by a correlated neural network.

\paragraph{Executable verification.} Unit tests, property-based tests, and simulation checks
provide selection under computational ground truth. Code that passes tests satisfies
constraints external to the generator. This may help explain why LLMs have become effective
at coding tasks: the code interpreter provides built-in external selection. The interpreter
does not care how confident the model was; the code runs or it throws an error.

\paragraph{Fresh context evaluation as an engineering heuristic.} The same model under
fresh context, without the reasoning chain that produced the candidate, may provide more
independent critique than same-context self-evaluation. We stress that fresh-context
same-model evaluation does \emph{not} satisfy the strict independence conditions of
Corollary~\ref{cor:external_selection}: shared weights and training distribution mean
parameter-level blind spots remain, so the conditional independence $S \perp T \mid (G, Z)$
is not broken in the required sense. Rather, context separation should be understood as a
\emph{pragmatic engineering heuristic} that partially reduces error correlation by removing
the generator's reasoning trace from the evaluator's context, without claiming to achieve
the theoretical standard. In the absence of a perfect solution, it serves as a low-cost
intervention that alters the activation context even when it cannot alter the underlying
weights. Whether this partial reduction is sufficient to yield meaningful practical benefit
is an empirical question addressed in planned future work
(Section~\ref{sec:observations}).

\citet{tsui2025selfcorrection} reports that even minimal context perturbation can
substantially affect performance: the ``Wait'' intervention, which injects a single
correction marker into the model's own reasoning trace, reduced the Self-Correction Blind
Spot by 89.3\%. This performance result is consistent with context disruption reducing
error correlation, but the evidence establishes the performance effect, not the formal
mechanism. Whether the intervention works by altering the shared failure structure $Z$,
by changing the evaluator's prediction target, or through some other mechanism is not
established by the cited data and remains a conjecture.

\begin{definition}[Context separation]
\label{def:context_sep}
Two evaluation contexts are \emph{separated} if the evaluating context has no access to:
(1) the generation trace (intermediate reasoning steps), (2) the prompt scaffolding
(instructions that shaped generation), or (3) hidden state from the generation process.
\end{definition}

\paragraph{Limitations of same-model context separation.} Fresh context reduces
context-level error correlation by removing the generator's intermediate reasoning trace and
local prompt scaffolding. However, it does not remove correlated failure modes that originate
in the model's parameters or training distribution. Same weights means shared inductive
biases remain. Context separation is therefore a partial solution: it breaks correlation
introduced \textit{during generation}, but not correlation baked into the model itself. For
maximum independence, context separation should be combined with model diversity or external
tools.

\paragraph{Numerical invariants.} Dimensional analysis, conservation laws, symmetry
constraints, and sanity bounds provide selection under physical ground truth. A derivation
that violates energy conservation fails regardless of how convincing it sounds.

\paragraph{Retrieval-grounded checking.} Citation verification against a fixed corpus, exact
quote attribution, and fact-checking against authoritative sources provide selection under
documentary ground truth.

\paragraph{Independent critics.} Models with different training data, different
architectures, or different optimization objectives have partially independent failure modes.

\subsection{Why Independence Matters}

Perfect independence is not required. What matters is that the joint failure probability is
lower than individual failure:
\begin{equation}
\Pr(E_{S_1} \cap E_{S_2} \mid E_G) < \Pr(E_{S_1} \mid E_G)
\end{equation}

Even partially independent selectors compound evidence. This explains empirically why
diverse multi-agent panels outperform single-agent self-evaluation even when individual
agents have similar capability~\citep{du2023debate, liang2023debate}.

\subsection{Persona-Based Diversity}

The predictive mechanism described in Section 3.6 suggests a strategy for achieving
decorrelated evaluation within a single model: cast the evaluator as different expert types.

When prompted as ``a rigorous mathematician checking for proof gaps,'' the model predicts
what such an expert would say. When prompted as ``a skeptical physicist checking dimensional
consistency,'' it predicts different behavior. When prompted as ``a journal referee looking
for reasons to reject,'' different still.

Each persona predicts a different distribution of expert behavior, with different priorities,
different blind spots, and different failure modes. An error that survives the mathematician
may not survive the physicist. A claim that passes the physicist may not pass the referee.

This creates epistemic diversity without model diversity. Multiple personas, each on clean
context, provide partially independent evaluations that can be aggregated. The decorrelation
arises not from different weights but from different prediction targets.

Implementation details (specific persona prompts, aggregation logic, consensus mechanisms)
are left to future work. Whether persona-based diversity achieves sufficient decorrelation
to satisfy Proposition~\ref{prop:multi_agent}'s requirements is an open empirical question.

%%%%%%%%%%%%%%%%%%%%%%%%%%%%%%%%%%%%%%%%%%%%%%%%%%%%%%%%%%%%%%%%%%%%%%%%%%%%%%%
\section{Context-Separated Architecture: Design Principles}
\label{sec:architecture}
%%%%%%%%%%%%%%%%%%%%%%%%%%%%%%%%%%%%%%%%%%%%%%%%%%%%%%%%%%%%%%%%%%%%%%%%%%%%%%%

The analysis above characterizes \emph{conditions} under which self-evaluation fails and
external selection helps. We now describe design principles for an architecture motivated by
that analysis. This section is a design proposal, not a validated system: we articulate what
properties a workflow should have to address correlated error, and how those properties can
be implemented. Whether implementations satisfying these properties perform as expected is an
empirical question addressed in planned future work (Section~\ref{sec:observations}).

The principles are organized around three functions: high-entropy generation to explore the
hypothesis space, external selection to filter candidates under criteria that are not
corrupted by the generator's blind spots, and feedback to update generation based on what
survives selection.

\subsection{Design Goals}

\begin{enumerate}
    \item \textbf{Maximize exploration}: Generate diverse candidates without premature filtering
    \item \textbf{Ensure rigor}: Select only candidates surviving external validation
    \item \textbf{Enable iteration}: Feed selection results back to improve generation
    \item \textbf{Preserve human judgment}: Surface candidates for human review, don't
    replace human decision-making
\end{enumerate}

\subsection{Architecture Components}

\begin{definition}[Generator]
The generator component $\mathcal{G}$ produces candidate hypotheses at high entropy,
exploring the space of possibilities without prejudice.
\end{definition}

\begin{definition}[Selector]
The selector component $\mathcal{S}$ evaluates candidates against external criteria at low
entropy, selecting configurations that survive validation.
\end{definition}

\begin{definition}[Feedback]
The Feedback component $\mathcal{B}$ updates generation context based on selection outcomes,
biasing future proposals toward surviving structures.
\end{definition}

\begin{equation}
\boxed{
\begin{aligned}
\mathcal{G}&: C \to \{h_1, \ldots, h_k\} && \text{(High-entropy generation)} \\
\mathcal{S}&: (h_i, C) \to (s_i, r_i) && \text{(External selection + rationale)} \\
\mathcal{B}&: \{(h_i, s_i, r_i)\} \to C' && \text{(Context update)}
\end{aligned}
}
\end{equation}

\subsection{High-Entropy Generation}

The generator component should:
\begin{itemize}
    \item Operate at high temperature or use explicit diversity objectives
    \item Generate candidates spanning the hypothesis space, including edge cases
    \item Avoid premature self-filtering that would narrow the search
    \item Include alternative assumptions and boundary conditions
\end{itemize}

Implementation options:
\begin{itemize}
    \item High-temperature sampling from a single model
    \item Ensemble sampling from multiple models
    \item Structured exploration of assumption variations
    \item Adversarial generation targeting unexplored regions
\end{itemize}

\subsection{External Selection}

The selector component should:
\begin{itemize}
    \item Operate at low temperature for consistency
    \item Apply explicit checklists and external tools
    \item Produce structured verdicts with rationales
    \item Flag uncertainty rather than forcing binary decisions
\end{itemize}

Selection criteria hierarchy:
\begin{enumerate}
    \item \textbf{Formal}: Does it compile/prove/type-check?
    \item \textbf{Executable}: Does it pass tests?
    \item \textbf{Numerical}: Does it satisfy invariants?
    \item \textbf{Grounded}: Do citations check out?
    \item \textbf{Adversarial}: Does it survive independent critique?
\end{enumerate}

\subsection{Learning from Selection}

The feedback component should:
\begin{itemize}
    \item Add constraints that killed candidates to future prompts
    \item Preserve successful patterns as templates
    \item Escalate ambiguous survivors to human review
    \item Track failure modes for architecture improvement
\end{itemize}

\subsection{Distinguishing from Related Architectures}

\begin{table}[h]
\centering
\caption{Comparison with Related Architectures}
\begin{tabular}{llll}
\toprule
\textbf{Architecture} & \textbf{Selection criterion} & \textbf{Goal} & \textbf{External?} \\
\midrule
GAN & Discriminator fooled & Realism & No (co-trained) \\
RLHF & Human preference & Alignment & Partially \\
Self-consistency & Agreement across samples & Confidence & No (correlated) \\
Context-separated & External validation & Truth & Yes \\
\bottomrule
\end{tabular}
\end{table}

The key distinction: context-separated evaluation uses persistence under external criteria
as the selection signal, not realism (GAN), preference (RLHF), or self-agreement
(self-consistency).

\subsection{Implementation Sketch}

\begin{verbatim}
GENERATE(context, diversity_target):
    candidates = []
    for i in 1..k:
        hypothesis = generate(context, temperature=HIGH)
        if diversity_check(hypothesis, candidates):
            candidates.append(hypothesis)
    return candidates

SELECT(hypothesis, context):
    verdict = {passed: True, checks: [], rationale: ""}

    # Formal checks
    if has_proof_component(hypothesis):
        proof_result = lean_check(hypothesis.proof)
        verdict.checks.append(("formal", proof_result))
        if not proof_result.success:
            verdict.passed = False

    # Numerical checks
    for invariant in domain_invariants:
        inv_result = check_invariant(hypothesis, invariant)
        verdict.checks.append(("numerical", inv_result))
        if not inv_result.success:
            verdict.passed = False

    # Adversarial checks
    for critic in independent_critics:
        critique = critic.evaluate(hypothesis)
        verdict.checks.append(("adversarial", critique))
        if critique.fatal_flaw:
            verdict.passed = False

    verdict.rationale = synthesize_rationale(verdict.checks)
    return verdict

FEEDBACK(candidates, verdicts, context):
    new_context = context
    for (h, v) in zip(candidates, verdicts):
        if not v.passed:
            new_context.add_constraint(v.rationale)
        else:
            new_context.add_template(h)
    new_context.survivors = [h for (h,v) in zip(...) if v.passed]
    return new_context

MAIN_LOOP(initial_context):
    context = initial_context
    all_survivors = []
    for round in 1..max_rounds:
        candidates = GENERATE(context, diversity_target)
        verdicts = [SELECT(h, context) for h in candidates]
        context = FEEDBACK(candidates, verdicts, context)
        all_survivors.extend(context.survivors)
        if convergence_criterion(context):
            break
    return prioritize_for_human_review(all_survivors)
\end{verbatim}

\subsection{Same-Model Implementation via Context Separation}

An important observation: this architecture does not require multiple models. The same model
under fresh context, without the reasoning chain that produced the candidate, can serve as
the selector. The key is context separation, not model separation.

\begin{verbatim}
CONTEXT_SEPARATED_EVALUATION(problem):
    # Context A: Generate with prediction
    context_a = fresh_context()
    response_a = model(context_a,
        "First predict what you think the answer is, then solve: " + problem)

    # Context B: Steelman AND attack (fresh, no access to reasoning)
    context_b = fresh_context()  # Critical: no shared state
    critique = model(context_b,
        "Here is a proposed solution. Provide both:
         1) Steelman: the strongest case FOR this solution
         2) Attack: the strongest case AGAINST this solution
         Solution: " + response_a.answer_only)

    # Context A: Judge coherence
    judgment = model(context_a,
        "A critic provided steelman and attack arguments.
         Which is more coherent with the problem structure?
         Steelman: " + critique.steelman +
         "Attack: " + critique.attack)

    return judgment
\end{verbatim}

This implementation has several conjectured advantages, pending empirical validation
(Section~\ref{sec:observations}):
\begin{itemize}
    \item \textbf{Favorable computational tradeoff}: Compared against extended
    chain-of-thought in a single context, fresh-context evaluation avoids the quadratic
    attention cost of very long contexts; the total token cost may be lower for problems
    requiring many reasoning steps. This comparison does not hold against a simple
    single-pass evaluation.
    \item \textbf{Reduced context-level correlation}: Context B cannot see Context A's
    reasoning trace, only the output. This removes one potential source of correlated error;
    parameter-level correlation from shared weights remains.
    \item \textbf{Simultaneous opposition}: Requesting both steelman and attack forces the
    model to genuinely consider both sides rather than anchoring on one.
    \item \textbf{Temperature control}: High temperature in generation (exploration), low
    temperature in critique (precision).
\end{itemize}

\subsection{Initial Observations and Future Validation}
\label{sec:observations}

This section describes qualitative observations that motivated the theoretical framework
above. These are not controlled experiments and we explicitly disclaim any empirical
conclusions from them. Our goal here is transparency about the origins of this work, not
empirical proof.

During development of this methodology, we observed the following patterns when submitting
manuscripts to fresh-context instances of commercial language models:

\begin{itemize}
    \item \textbf{Prompt sensitivity}: Minor framing changes (e.g., ``thoughts'' vs
    ``honest review'') produced noticeably different evaluations of identical content.
    \item \textbf{In-context persuasion}: Critics who heard the author's defense sometimes
    revised harsh assessments to positive ones, suggesting context sharing may correlate
    evaluator judgment with author framing.
    \item \textbf{Fresh-context disagreement}: Multiple fresh-context evaluations of the
    same manuscript sometimes disagreed with each other, while self-evaluation produced more
    consistent (but potentially unreliable) agreement.
\end{itemize}

These observations are anecdotal and indicative only. They motivated the formal analysis in
Sections 2--3 but do not constitute validation of it. No protocols, sampling schemes, or
statistical analyses were applied.

\paragraph{Falsifiability and future work.} This analysis makes testable predictions:
(1) self-evaluation should show higher error correlation than context-separated evaluation;
(2) fresh-context critics should catch errors that in-context self-critique misses;
(3) disagreement among independent evaluators should correlate with actual uncertainty
about correctness.

Rigorous validation requires controlled experiments across multiple models, systematic
variation of prompts and system configurations, and proper statistical methodology. This is
beyond the scope of a methodology paper and is left to future work. A companion paper will
detail empirical results across model families, prompt variations, and evaluation criteria.

\paragraph{Collaboration invited.} We recognize that experimental design for evaluating LLM
self-assessment is methodologically challenging. We welcome collaboration on reducing bias
in experimental protocols. Complete conversation logs from our preliminary observations are
available on request.

\paragraph{Additional observations.} We also observed evaluation sensitivity to surface
features: formatting changes, word choice (e.g., ``novel'' vs ``improved''), and
acknowledged prestige bias when we asked models directly whether author reputation would
affect their assessment. These patterns are consistent with LLM evaluation inheriting human
biases from training data, but we do not claim these as experimental findings.

\paragraph{Negative results.} Context separation does not eliminate the evaluation format
described above. When evaluating low-quality work, the balanced positive/negative structure
persists; the model fills the expected ``positive points'' slots with increasingly tenuous
or hallucinated content rather than breaking format to deliver an unbalanced negative
assessment. Context separation reduces correlated error but does not override the formatting
prior.

\paragraph{Stop condition.} In practice, iterative critique terminates when critics begin
producing hallucinated criticism: attacks referencing problems not present, repeating
addressed points, or focusing on irrelevant details. Detection requires human judgment.

%%%%%%%%%%%%%%%%%%%%%%%%%%%%%%%%%%%%%%%%%%%%%%%%%%%%%%%%%%%%%%%%%%%%%%%%%%%%%%%
\section{Threat Model and Mitigations}
%%%%%%%%%%%%%%%%%%%%%%%%%%%%%%%%%%%%%%%%%%%%%%%%%%%%%%%%%%%%%%%%%%%%%%%%%%%%%%%

To clarify the architecture's robustness, we enumerate failure modes and mitigations.

\subsection{Attack Surfaces}

\paragraph{Correlated critics.} If critic models share training data, architecture, or
optimization objectives with the generator, error correlation persists despite apparent
multi-agent structure. This is the primary failure mode the architecture addresses.

\paragraph{Prompt injection on selection.} Adversarial inputs could manipulate critic
behavior, causing systematic acceptance of flawed candidates.

\paragraph{Spoofable external checks.} If ``external'' verification tools can be fooled
(e.g., tests that don't actually test the claimed property), the selection signal degrades.

\paragraph{Human confirmation bias.} The human operator may preferentially accept candidates
that confirm prior beliefs, reintroducing correlated error at the final selection stage.

\paragraph{Feedback gaming.} If the feedback mechanism is predictable, generators could
learn to produce candidates that pass selection without satisfying underlying criteria.

\subsection{Mitigations}

\begin{table}[h]
\centering
\caption{Threat Mitigations}
\begin{tabular}{ll}
\toprule
\textbf{Threat} & \textbf{Mitigation} \\
\midrule
Correlated critics & Model diversity (different families, training data, architectures) \\
Prompt injection & Structured verdict formats, input sanitization, separate contexts \\
Spoofable checks & Hard external criteria (formal proofs, physical measurements) \\
Human bias & Adversarial critics, blind review protocols, explicit checklists \\
Feedback gaming & Diverse selection criteria, periodic architecture audits \\
\bottomrule
\end{tabular}
\end{table}

No mitigation is complete. The goal is defense in depth: multiple independent failure modes
such that exploiting one does not compromise the entire system.

%%%%%%%%%%%%%%%%%%%%%%%%%%%%%%%%%%%%%%%%%%%%%%%%%%%%%%%%%%%%%%%%%%%%%%%%%%%%%%%
\section{Implications}
%%%%%%%%%%%%%%%%%%%%%%%%%%%%%%%%%%%%%%%%%%%%%%%%%%%%%%%%%%%%%%%%%%%%%%%%%%%%%%%

\emph{Note on scope: The implications below are motivated by the formal analysis but are
not logically derived from it. They represent engineering and design suggestions consistent
with the theory's conditional conclusions, not predictions guaranteed by the theorems.}

\subsection{For AI-Assisted Workflows}

The formal analysis identifies conditions under which scaling a single model's self-evaluation
may be insufficient for reliable validation in generate-then-judge workflows. Where those
conditions apply, the bottleneck is not generation capability but validation under criteria
that are genuinely external to the generator's failure structure. Investment in external
selection infrastructure (formal verification tools, test generation, invariant libraries,
diverse critic ensembles) may yield returns complementary to larger single models in such
settings.

\subsection{For Reduced Human Oversight}

The analysis identifies conditions under which workflows with reduced human oversight may
require selection mechanisms with low correlation to generator error. Where those conditions
apply, current approaches using self-critique or same-family judge models may not provide
the independence required. Reducing human oversight while maintaining reliability in such
settings may require:
\begin{itemize}
    \item Formal verification covering the relevant claim space
    \item Executable tests providing ground truth
    \item Evaluation architectures with independently constrained failure modes
\end{itemize}

In settings where these mechanisms are unavailable, human review remains an important
external criterion. Whether a given deployment is in the regime the theory analyzes is
itself an empirical question.

\subsection{For Human-AI Collaboration}

The architecture's goal is not to replace human judgment but to concentrate it. By filtering
candidates through external selection, the system surfaces hypotheses worthy of human
attention. This addresses a scalability consideration: humans cannot evaluate 200 hypotheses,
but they can evaluate 2--3 survivors of rigorous automated filtering.

\subsection{Deployment Considerations for Research Applications}

Researchers needing adversarial evaluation may benefit from deployments with full prompt
control. Commercial APIs constrain the system prompt, layering user instructions atop a
fixed ``helpful assistant'' foundation. Local deployment with customizable system prompts
enables research-specific evaluation personas (e.g., ``Skeptic,'' ``Hostile Reviewer'').
This is not a criticism of commercial design choices, which appropriately prioritize safety
and broad utility, but an observation about specialization requirements for research workflows.

\subsection{Open Questions}

\begin{enumerate}
    \item What constraints on $Z$ are both empirically tractable and sufficient to make
    the bounds of Section~\ref{sec:self_eval_fails} non-vacuous?
    \item Does fresh-context same-model evaluation yield measurable reduction in error
    correlation in practice, even though it does not satisfy the formal independence
    conditions?
    \item Can formal verification scale to cover more scientific domains?
    \item What is the minimum selector diversity required for the conditions of
    Proposition~\ref{prop:multi_agent} to hold in practice?
    \item Can selection criteria themselves be learned without introducing correlation with
    the generator's failure structure?
\end{enumerate}

\subsection{Future Work}

The preliminary observations in Section~\ref{sec:observations} motivated this theoretical
work but do not constitute rigorous validation. The most important open problems, in order
of priority, are:
\begin{itemize}
    \item \textbf{Operationalizing $Z$ using existing taxonomies}: Apply the ErrorAtlas
    17-category taxonomy~\citep{erroratlas2026} or the Song et al.\ two-axis
    framework~\citep{song2026reasoning} as candidate a priori $Z$ specifications, then
    test whether the conditional independence $S \perp T \mid (G, Z)$ holds empirically
    using existing benchmark data stratified by error category
    \item \textbf{Measuring $\Pr(A=1 \mid T=0)$ using existing infrastructure}:
    CriticBench~\citep{lin2024criticbench} and the P(True) framework~\citep{kadavath2022ptrue}
    already measure binary correctness classification against known ground truth. Running
    these protocols in self-evaluation mode (same model generates and critiques) vs.\
    external-evaluation mode would directly test the proposition's conditions and measure
    whether deployed systems are in the failure regime
    \item \textbf{Testing the joint independence assumption}: Using error-stratified
    benchmarks to measure whether multiple self-evaluation rounds provide incrementally
    less information than expected under independence, as predicted by
    Lemma~\ref{lem:repeated}
    \item \textbf{Measuring context separation's practical effect}: Controlled experiments
    comparing self-evaluation to context-separated evaluation, with proper measurement of
    error correlation using inter-rater reliability metrics
\end{itemize}
Until such validation is complete, we explicitly disclaim strong empirical conclusions.
The contribution of this paper is the theoretical framework and the identification of
concrete empirical pathways to testing it; empirical confirmation remains future work.

%%%%%%%%%%%%%%%%%%%%%%%%%%%%%%%%%%%%%%%%%%%%%%%%%%%%%%%%%%%%%%%%%%%%%%%%%%%%%%%
\section{Related Work}
%%%%%%%%%%%%%%%%%%%%%%%%%%%%%%%%%%%%%%%%%%%%%%%%%%%%%%%%%%%%%%%%%%%%%%%%%%%%%%%

\paragraph{LLM self-correction: the empirical foundation.} \citet{huang2024selfcorrect}
established empirically that ``large language models cannot self-correct reasoning yet,''
showing that without external feedback, self-correction attempts often fail or degrade
performance. Our work offers one possible explanation: correlated error between generator and
evaluator can render self-evaluation non-identifying. Their finding that multi-agent critique
with same-model copies performed ``no better than self-consistency'' is consistent with our
analysis: identical models share error distributions, so adding copies may not break
correlation. Our analysis suggests a fix: context separation or genuinely independent
evaluation channels.

\paragraph{The Self-Correction Blind Spot.} \citet{tsui2025selfcorrection} provides
important empirical context for the framework. The study demonstrates that LLMs can correct
identical errors when presented as external input but fail to correct those same errors in
their own outputs, a phenomenon termed the Self-Correction Blind Spot. An average 64.5\%
failure rate is measured across 14 models on three purpose-built benchmarks (SCLI5,
GSM8K-SC, PRM800K-SC). The ``Wait'' intervention, which injects correction markers into
the model's reasoning trace, reduced the blind spot by 89.3\%, activating latent
capabilities without changing model weights. We treat these findings as empirical context
rather than formal validation: the mapping between ``failure to self-correct'' and the
theorem's formal quantities ($\Pr(A=1 \mid T=0)$, conditional independence relations) is
not established, and the mechanism by which ``Wait'' produces its effect is not identified
by the performance evidence. The findings are consistent with the shared-blind-spot model
but do not confirm it. Where we provide theoretical conditions under which self-evaluation
may fail, Tsui provides empirical measurement of failure rates and a candidate intervention
whose mechanism remains to be characterized.

\paragraph{Self-consistency and majority voting.} \citet{wang2023selfconsistency} demonstrated
that sampling multiple reasoning paths and taking the majority answer improves accuracy. This
works when errors are uncorrelated across samples. Our analysis clarifies the limitation: if
all samples share the same systematic bias (high error correlation), majority voting cannot
help. The gains from self-consistency diminish as correlation increases, explaining why the
technique works better for some tasks than others.

\paragraph{Multi-agent debate and verification.} Multi-agent systems including
AutoGen~\citep{wu2023autogen}, CAMEL~\citep{li2023camel}, MetaGPT~\citep{hong2023metagpt},
and debate frameworks~\citep{du2023debate, liang2023debate} explore collaborative reasoning.
\citet{chen2024reconcile} found that model \textit{diversity} among agents was important for
performance gains, consistent with the idea that breaking error correlation matters. We
attempt to formalize why multi-agent verification can help (decorrelated failure modes) and
suggest that context separation within a single model may achieve similar benefits.

\paragraph{Process supervision.} \citet{lightman2023process} showed that supervising each
reasoning step (process supervision) significantly outperforms supervising only final answers
(outcome supervision). This aligns with our analysis: per-step external feedback breaks the
model's ``solo reasoning bubble,'' preventing error accumulation within a single correlated
context. Process supervision is an instance of external selection applied during training.

\paragraph{External verification and tool use.} Training separate verifier
models~\citep{cobbe2021verifier} achieved large gains on math problems. Integration with
formal provers~\citep{yang2023leandojo, polu2020generative, baba2025prover} provides external
selection via mathematical ground truth. The CRITIC framework~\citep{gou2023critic} showed
that tool-interactive critiquing improves self-correction. These results are consistent with
the view that ``external'' can mean different models, formal tools, or execution environments.

\paragraph{Apparent counterexamples.} Some work reports successful self-correction:
Self-Refine~\citep{madaan2023selfrefine} for iterative text improvement,
Reflexion~\citep{shinn2023reflexion} for agent learning, and Constitutional
AI~\citep{bai2022constitutional} for safety. We reconcile these with our analysis by noting
they typically address (a) style and format rather than deep reasoning, (b) cases where
oracle feedback is implicitly available, or (c) safety constraints that are well-represented
in training data. When \citet{huang2024selfcorrect} removed oracle feedback from
self-correction setups, improvements vanished, suggesting that apparent self-correction may
often rely on hidden external signals.

\paragraph{Correlated errors and LLM-as-a-Judge biases.}
\citet{kim2025correlated} provide the most direct large-scale measurement of the correlated
error assumption underlying this framework. Across 350+ LLMs on two leaderboards, same-family
models agree on the same wrong answer approximately 60\% of the time; critically, error
correlation increases with model capability. This directly instantiates the high-$\kappa$
regime the framework analyzes, and the authors demonstrate that LLM-as-judge validity is
undermined by these correlations. \citet{panickssery2024selfpreference} found significant
self-preference bias: LLMs assign higher scores to outputs with lower perplexity.
\citet{li2024llmjudgebias} identified 12 major latent biases in LLM-as-a-Judge systems.
These empirical findings align with our theoretical analysis: evaluation that shares the
generator's distribution will exhibit correlated error.

\paragraph{Ensemble diversity and error decorrelation.} The insight that independent errors
enable reliable aggregation is foundational in ensemble learning~\citep{krogh1995ensemble}.
We apply this insight to LLM self-evaluation, where error correlation may be particularly
relevant due to shared training data, weights, and context. We connect ensemble theory to
information theory: conditional mutual information $I(T; S \mid G)$ can approach zero under
high correlation~\citep{cover2006elements}, which would explain why self-evaluation becomes
uninformative in such cases.

\paragraph{AI for science.} AlphaFold~\citep{alphafold} demonstrates AI solving well-defined
scientific problems with clear evaluation criteria. Our focus is the less-structured setting
of hypothesis generation and validation, where ground truth is not known in advance and
external selection must be actively constructed.

%%%%%%%%%%%%%%%%%%%%%%%%%%%%%%%%%%%%%%%%%%%%%%%%%%%%%%%%%%%%%%%%%%%%%%%%%%%%%%%
\section{Conclusion}
%%%%%%%%%%%%%%%%%%%%%%%%%%%%%%%%%%%%%%%%%%%%%%%%%%%%%%%%%%%%%%%%%%%%%%%%%%%%%%%

The paper's primary result is the repeated self-critique bound: under joint conditional
independence of evaluations given the shared failure structure, $k$ rounds of self-critique
provide information about correctness bounded by what the shared latent variable $Z$
mediates---not by any independent channel. Confidence accumulated through repeated
self-evaluation therefore reflects the shared failure structure rather than independently
accumulated evidence. Whether the joint conditional independence assumption holds for real
LLM systems is an empirically testable question, not a consequence of the formalism.

The multi-agent advantage proposition gives the constructive counterpart: a selector
provides positive evidence about correctness if and only if it achieves a bounded
false-acceptance rate \emph{and} a true-acceptance rate exceeding that bound. Both
quantities are directly measurable for real systems. This gives the diagnostic framework
its practical content: instead of asking whether self-evaluation ``fails'' in some abstract
sense, practitioners can measure whether deployed selectors satisfy these conditions.

The information-theoretic bound (Theorem~\ref{thm:info_bound}) is best understood as
scaffolding for these results. It shows that under a shared-blind-spot modeling assumption,
self-evaluation is bounded in what it can add --- but the bound is vacuous without an
independently constrained latent variable $Z$. We foreground this limitation rather than
minimizing it. Constructing a falsifiable, a priori characterization of $Z$ for real LLM
systems is the key open problem the framework identifies, and we propose concrete approaches
in the future work section.

The design principles we describe are motivated by the analysis but not derived from it.
Same-model context separation is an engineering heuristic --- it removes the generation
trace from the evaluation context without breaking the parameter-level correlations the
theory identifies as the source of difficulty. We present it as a practical starting point,
with the explicit prediction that it will provide partial benefit, and defer validation of
that prediction to empirical future work.

Recent findings by \citet{tsui2025selfcorrection} --- that models fail to self-correct
at high rates, and that minimal context disruption substantially reduces blind-spot rates
--- are cited as empirical context, not formal validation. The mapping between those
findings and the theorem's formal quantities is not established; confirming that real
systems are in the failure regime requires controlled experiments with explicit acceptance
decisions, which we leave to future work.

The framework's value is diagnostic: it identifies what to measure to determine whether a
system is in the failure regime, and what properties an external selector must have to
escape it. The bottleneck in reliable AI-assisted workflows is often validation rather than
generation. The framework makes that bottleneck precise.

%%%%%%%%%%%%%%%%%%%%%%%%%%%%%%%%%%%%%%%%%%%%%%%%%%%%%%%%%%%%%%%%%%%%%%%%%%%%%%%
\section*{Use of AI Tools}
%%%%%%%%%%%%%%%%%%%%%%%%%%%%%%%%%%%%%%%%%%%%%%%%%%%%%%%%%%%%%%%%%%%%%%%%%%%%%%%

AI tools assisted with drafting and editing. All ideas, methodology, theorems, and technical
content are the author's own work.

%%%%%%%%%%%%%%%%%%%%%%%%%%%%%%%%%%%%%%%%%%%%%%%%%%%%%%%%%%%%%%%%%%%%%%%%%%%%%%%
\section*{Acknowledgments}
%%%%%%%%%%%%%%%%%%%%%%%%%%%%%%%%%%%%%%%%%%%%%%%%%%%%%%%%%%%%%%%%%%%%%%%%%%%%%%%

The author thanks anonymous reviewers for helpful feedback.

%%%%%%%%%%%%%%%%%%%%%%%%%%%%%%%%%%%%%%%%%%%%%%%%%%%%%%%%%%%%%%%%%%%%%%%%%%%%%%%
\section*{Declarations}
%%%%%%%%%%%%%%%%%%%%%%%%%%%%%%%%%%%%%%%%%%%%%%%%%%%%%%%%%%%%%%%%%%%%%%%%%%%%%%%

\textbf{Funding:} This research received no external funding.

\textbf{Conflicts of Interest:} The author declares no conflicts of interest.

\bibliographystyle{plainnat}
\begin{thebibliography}{99}

\bibitem[Baba et al.(2025)]{baba2025prover}
Kaito Baba, Chaoran Liu, Shuhei Kurita, and Akiyoshi Sannai.
\newblock Prover Agent: An agent-based framework for formal mathematical proofs.
\newblock \emph{arXiv preprint arXiv:2506.19923}, 2025.

\bibitem[Bai et al.(2022)]{bai2022constitutional}
Yuntao Bai et al.
\newblock Constitutional {AI}: Harmlessness from {AI} feedback.
\newblock \emph{arXiv preprint arXiv:2212.08073}, 2022.

\bibitem[Bender et al.(2021)]{bender2021parrots}
Emily~M. Bender, Timnit Gebru, Angelina McMillan-Major, and Shmargaret Shmitchell.
\newblock On the Dangers of Stochastic Parrots: Can Language Models Be Too Big?
\newblock \emph{Proceedings of the 2021 ACM Conference on Fairness, Accountability, and
Transparency (FAccT)}, pages 610--623, 2021.

\bibitem[Chen et al.(2024)]{chen2024reconcile}
Justin Chih-Yao Chen, Swarnadeep Saha, and Mohit Bansal.
\newblock ReConcile: Round-table conference improves reasoning via consensus among diverse
{LLMs}.
\newblock \emph{Proceedings of the 62nd Annual Meeting of the Association for Computational
Linguistics (ACL)}, 2024.

\bibitem[Cobbe et al.(2021)]{cobbe2021verifier}
Karl Cobbe et al.
\newblock Training verifiers to solve math word problems.
\newblock \emph{arXiv preprint arXiv:2110.14168}, 2021.

\bibitem[Cover and Thomas(2006)]{cover2006elements}
Thomas M.\ Cover and Joy A.\ Thomas.
\newblock \emph{Elements of Information Theory}.
\newblock Wiley-Interscience, 2nd edition, 2006.

\bibitem[Du et al.(2024)]{du2023debate}
Yilun Du, Shuang Li, Antonio Torralba, Joshua B.\ Tenenbaum, and Igor Mordatch.
\newblock Improving factuality and reasoning in language models through multiagent debate.
\newblock \emph{Proceedings of the 41st International Conference on Machine Learning (ICML)},
2024.

\bibitem[G\"{o}del(1931)]{godel1931}
Kurt G\"{o}del.
\newblock \"{U}ber formal unentscheidbare S\"{a}tze der Principia Mathematica und verwandter
Systeme I.
\newblock \emph{Monatshefte f\"{u}r Mathematik und Physik}, 38:173--198, 1931.

\bibitem[Gou et al.(2023)]{gou2023critic}
Zhibin Gou et al.
\newblock {CRITIC}: Large language models can self-correct with tool-interactive critiquing.
\newblock \emph{arXiv preprint arXiv:2305.11738}, 2023.

\bibitem[Hong et al.(2023)]{hong2023metagpt}
Sirui Hong et al.
\newblock MetaGPT: Meta programming for a multi-agent collaborative framework.
\newblock \emph{arXiv preprint arXiv:2308.00352}, 2023.

\bibitem[Huang et al.(2024)]{huang2024selfcorrect}
Jie Huang, Xinyun Chen, Swaroop Mishra, Huaixiu Steven Zheng, Adams Wei Yu, Xinying Song,
and Denny Zhou.
\newblock Large language models cannot self-correct reasoning yet.
\newblock \emph{Proceedings of the Twelfth International Conference on Learning
Representations (ICLR)}, 2024.

\bibitem[Jumper et al.(2021)]{alphafold}
John Jumper et al.
\newblock Highly accurate protein structure prediction with AlphaFold.
\newblock \emph{Nature}, 596:583--589, 2021.

\bibitem[Kong et al.(2023)]{kong2023roleplay}
Aobo Kong, Shiwan Zhao, Hao Chen, Qicheng Li, Yong Qin, Ruiqi Sun, and Xin Zhou.
\newblock Better Zero-Shot Reasoning with Role-Play Prompting.
\newblock \emph{Proceedings of the 2024 Conference of the North American Chapter of the
Association for Computational Linguistics (NAACL)}, 2024.

\bibitem[Krogh and Vedelsby(1995)]{krogh1995ensemble}
Anders Krogh and Jesper Vedelsby.
\newblock Neural network ensembles, cross validation, and active learning.
\newblock \emph{Advances in Neural Information Processing Systems (NIPS)}, 1995.

\bibitem[Li et al.(2024)]{li2024llmjudgebias}
Jiayi Li et al.
\newblock Justice or prejudice? Quantifying biases in LLM-as-a-Judge.
\newblock \emph{arXiv preprint arXiv:2410.02736}, 2024.

\bibitem[Li et al.(2023)]{li2023camel}
Guohao Li et al.
\newblock CAMEL: Communicative agents for ``mind'' exploration of large language model society.
\newblock \emph{Advances in Neural Information Processing Systems}, 36, 2023.

\bibitem[Liang et al.(2023)]{liang2023debate}
Tian Liang et al.
\newblock Encouraging divergent thinking in large language models through multi-agent debate.
\newblock \emph{arXiv preprint arXiv:2305.19118}, 2023.

\bibitem[Lightman et al.(2023)]{lightman2023process}
Hunter Lightman et al.
\newblock Let's verify step by step.
\newblock \emph{arXiv preprint arXiv:2305.20050}, 2023.
\newblock (OpenAI; ICLR 2024).

\bibitem[Madaan et al.(2023)]{madaan2023selfrefine}
Aman Madaan et al.
\newblock Self-Refine: Iterative refinement with self-feedback.
\newblock \emph{Advances in Neural Information Processing Systems (NeurIPS)}, 2023.

\bibitem[Panickssery et al.(2024)]{panickssery2024selfpreference}
Arjun Panickssery et al.
\newblock Self-preference bias in LLM-as-a-Judge.
\newblock \emph{arXiv preprint arXiv:2410.21819}, 2024.

\bibitem[Polu and Sutskever(2020)]{polu2020generative}
Stanislas Polu and Ilya Sutskever.
\newblock Generative language modeling for automated theorem proving.
\newblock \emph{arXiv preprint arXiv:2009.03393}, 2020.

\bibitem[Pronin et~al.(2002)]{pronin2002bias}
Emily Pronin, Daniel~Y. Lin, and Lee Ross.
\newblock The bias blind spot: Perceptions of bias in self versus others.
\newblock \emph{Personality and Social Psychology Bulletin}, 28(3):369--381, 2002.

\bibitem[Sclar et al.(2024)]{sclar2024formatting}
Melanie Sclar, Yejin Choi, Yulia Tsvetkov, and Alane Suhr.
\newblock Quantifying Language Models' Sensitivity to Spurious Features in Prompt Design or:
How I Learned to Start Worrying about Prompt Formatting.
\newblock \emph{Proceedings of the Twelfth International Conference on Learning
Representations (ICLR)}, 2024.

\bibitem[Sharma et al.(2024)]{sharma2024sycophancy}
Mrinank Sharma, Meg Tong, Tomasz Korbak, David Duvenaud, Amanda Askell, Samuel~R. Bowman,
and Ethan Perez.
\newblock Towards Understanding Sycophancy in Language Models.
\newblock \emph{Proceedings of the Twelfth International Conference on Learning
Representations (ICLR)}, 2024.

\bibitem[Shinn et al.(2023)]{shinn2023reflexion}
Noah Shinn, Federico Cassano, Ashwin Gopinath, Karthik Narasimhan, and Shunyu Yao.
\newblock Reflexion: Language agents with verbal reinforcement learning.
\newblock \emph{Advances in Neural Information Processing Systems (NeurIPS)}, 2023.

\bibitem[Tsui(2025)]{tsui2025selfcorrection}
Ken Tsui.
\newblock Self-Correction Bench: Uncovering and Addressing the Self-Correction Blind Spot in
Large Language Models.
\newblock \emph{arXiv preprint arXiv:2507.02778}, 2025.

\bibitem[Tyen et al.(2024)]{tyen2024llmsstepbystep}
Gladys Tyen, Hassan Mansoor, Victor Carbune, Peter Chen, and Tony Mak.
\newblock LLMs Cannot Find Reasoning Errors, but Can Correct Them Given the Error Location.
\newblock \emph{Findings of the Association for Computational Linguistics: ACL 2024},
pages 13894--13908, 2024.

\bibitem[Wang et al.(2023)]{wang2023selfconsistency}
Xuezhi Wang et al.
\newblock Self-consistency improves chain of thought reasoning in language models.
\newblock \emph{Proceedings of the Eleventh International Conference on Learning
Representations (ICLR)}, 2023.

\bibitem[Wu et al.(2023)]{wu2023autogen}
Qingyun Wu et al.
\newblock AutoGen: Enabling next-gen LLM applications via multi-agent conversation.
\newblock \emph{arXiv preprint arXiv:2308.08155}, 2023.

\bibitem[Yang et al.(2023)]{yang2023leandojo}
Kaiyu Yang et al.
\newblock LeanDojo: Theorem proving with retrieval-augmented language models.
\newblock \emph{Advances in Neural Information Processing Systems}, 36, 2023.

\bibitem[Zheng et al.(2023)]{zheng2023helpful}
Mingqian Zheng, Jiaxin Pei, and David Jurgens.
\newblock Is ``A Helpful Assistant'' the Best Role for Large Language Models? A Systematic
Evaluation of Social Roles in System Prompts.
\newblock \emph{arXiv preprint arXiv:2311.10054}, 2023.

\bibitem[Zhuo et al.(2024)]{zhuo2024prosa}
Jingming Zhuo, Songyang Zhang, Xinyu Fang, Haodong Duan, Dahua Lin, and Kai Chen.
\newblock ProSA: Assessing and Understanding the Prompt Sensitivity of LLMs.
\newblock \emph{Findings of the Association for Computational Linguistics: EMNLP 2024}, 2024.

\bibitem[Ashury-Tahan et al.(2026)]{erroratlas2026}
Ofir Ashury-Tahan, Matan Mai, Elron Bandel, Michal Shmueli-Scheuer, and Leshem Choshen.
\newblock {ErrorMap} and {ErrorAtlas}: Charting the Failure Landscape of Large Language
Models.
\newblock \emph{arXiv preprint arXiv:2601.15812}, 2026.

\bibitem[Kadavath et al.(2022)]{kadavath2022ptrue}
Saurav Kadavath et al.
\newblock Language Models (Mostly) Know What They Know.
\newblock \emph{arXiv preprint arXiv:2207.05221}, 2022.

\bibitem[Kim et al.(2025)]{kim2025correlated}
Gwanghee Kim, Tanmay Garg, Vivek Peng, and Neha Garg.
\newblock Correlated Errors in Large Language Models.
\newblock \emph{Proceedings of the 42nd International Conference on Machine Learning
(ICML)}, 2025.

\bibitem[Lin et al.(2024)]{lin2024criticbench}
Zicheng Lin, Zhibin Gou, Tian Liang, Ruilin Luo, Haowei Liu, and Yujiu Yang.
\newblock {CriticBench}: Benchmarking {LLMs} for Critique-Correct Reasoning.
\newblock \emph{Findings of the Association for Computational Linguistics: ACL 2024}, 2024.

\bibitem[Song et al.(2026)]{song2026reasoning}
Yuxuan Song, Noah Han, and Noah Goodman.
\newblock Large Language Model Reasoning Failures.
\newblock \emph{Transactions on Machine Learning Research}, 2026.

\end{thebibliography}

\end{document}